\documentclass[11pt,a4paper]{article}
\usepackage[margin=25mm]{geometry}
\usepackage{amsmath,amssymb,amsthm,booktabs,hyperref}
\newtheorem{proposition}{Proposition}
\title{Chirality detected by the order of one word}
\author{Human Compression project}\date{}
\begin{document}\maketitle
We give finite certificates for twelve sporadic cases of the existing
chiral-polyhedron existence theorem of Leemans and Liebeck~\cite{LL}.
The same short word detects chirality in every case. The two Mathieu
pairs printed in that paper are retained verbatim; the other pairs are
explicit new witnesses to already known existence assertions. Six of
their ambient representations are acquired from our earlier, separately
reviewed certificates. None of these cases is claimed as a new
existence theorem or as a solution to an open conjecture.

\section{The common criterion}
If $w$ is a word in two letters, write $w(x,t)$ for its value at $x,t$.
\begin{proposition}\label{word}
Suppose $G=\langle x,t\rangle$ is a finite group, $t$ has order two,
and some word $w$ satisfies
\[
 |w(x,t)|\ne |w(x^{-1},t)|.
\]
No automorphism of $G$ can invert $x$ and fix $t$.
\end{proposition}
\begin{proof}
An automorphism with this action would send $w(x,t)$ to
$w(x^{-1},t)$. An isomorphism preserves element orders, contradicting
the displayed inequality. This excludes all automorphisms with the
stated action, including the involutory ones in the geometric criterion.
\end{proof}
The proposition is an elementary use of word invariants, not a new
automorphism theorem. It is useful because a single checked inequality
replaces a search through the automorphism group.

For completeness, suppose also that $Z(G)=1$, and put $y=x^{-1}t$.
An element of $\langle x\rangle\cap\langle y\rangle$ commutes with
both generators $x,y$, and hence belongs to $Z(G)$. Thus
\[
 \langle x\rangle\cap\langle y\rangle=1,\qquad (xy)^2=1.
\]
In all the examples below $|x|,|y|>2$. These are precisely the rank-three
rotation intersection and product conditions. We retain the geometric
reconstruction theorem of Schulte and Weiss~\cite{SW}: these data give
a directly regular or chiral polyhedron, and the automorphism in
Proposition~\ref{word} is the regularity alternative. Its absence gives
a chiral polyhedron of type $\{|x|,|y|\}$ with full automorphism group
$G$. The sporadic groups below are nonabelian simple groups, so their
centres are trivial. Their identification and simplicity are explicit
retained classification inputs.

\section{Twelve exact instances}
In every row use the single word
\[
 w(x,t)=x t x^2 t x^3 t.
\]
The last two columns differ in every case.
\begin{center}
\begin{tabular}{lrrrrr}\toprule
$G$&$|G|$&$|x|$&$|x^{-1}t|$&$|w(x,t)|$&$|w(x^{-1},t)|$\\\midrule
$M_{11}$&7920&11&8&8&3\\
$M_{12}$&95040&8&8&10&11\\
$M_{22}$&443520&8&7&7&5\\
$M_{23}$&10200960&11&15&14&6\\
$M_{24}$&244823040&23&15&15&12\\
$J_1$&175560&10&6&6&19\\
$J_2$&604800&15&15&10&15\\
$J_3$&50232960&10&9&4&17\\
$HS$&44352000&8&7&7&11\\
$He$&4030387200&14&15&17&21\\
$Suz$&448345497600&10&9&12&8\\
$McL$&898128000&30&8&14&11\\\bottomrule
\end{tabular}
\end{center}
These rows assert existence. They do not recover the source's ancillary
counts of nonisomorphic polyhedra or require its particular example type
when we supply a different witness.

The files \texttt{mathieu\_pairs.g}, \texttt{rank4\_pairs.g} and
\texttt{small\_pairs.g} specify every permutation. For $M_{23},M_{24}$,
the source gives $x,y$ and we set $t=xy$. The checker tests the two
permutations literally and computes their orders. For $M_{23}$ it
checks their actual membership in the named Mathieu group. For
$M_{24}$ the displayed labelling differs from the installed standard
one, so the checker constructs a bijective homomorphism onto that
named group. Equality of orders alone is not used to identify it.

For $M_{12},J_2,HS,He,Suz,J_3$, let $(r_1,r_2,r_3,r_4)$ be the acquired
four-involution marking. The saved list of thirty-two indices
$(i_1,\ldots,i_{32})$ defines
\[
          x=r_{i_1}\cdots r_{i_{32}},\qquad t=r_1.
\]
The whole marking, its ambient generators and the named-group
identification are bound inputs. For the first five groups the
ambient extension is, respectively,
\texttt{PrimitiveGroup(144,4)}, \texttt{PrimitiveGroup(100,2)},
\texttt{PrimitiveGroup(100,4)}, \texttt{PrimitiveGroup(2058,2)},
and \texttt{PrimitiveGroup(1782,2)}. Its derived subgroup is the
specified simple group. The $J_3:2$ representation and defining
presentation are acquired from the earlier certificate and its
Atlas inputs~\cite{Atlas}. We retain those named identifications;
we do not charge them as if the group name followed from an order.

The remaining four pairs are literal permutations in
\texttt{MathieuGroup(11)}, \texttt{MathieuGroup(22)},
the derived subgroup of \texttt{PrimitiveGroup(275,2)}, and
\texttt{PrimitiveGroup(266,1)}. The last catalogue record explicitly
identifies $J_1$; that identification is a retained input too.
The checker tests membership of both permutations in each named
group before comparing the generated subgroup's order with its order.

\section{Certificate and scope}
For every saved pair, the finite checker verifies $t^2=1$, $t\ne1$,
membership in the specified group, generation of the full group,
and each of the four orders in the table. Since the group is finite,
a contained subgroup of equal order is the whole group. This proves
generation from the saved data. Proposition~\ref{word} then proves
the absence of reversal. The geometric conclusion uses only the
retained reconstruction theorem and the already identified simple
group.

The witness searches are construction aids and are not replay inputs.
No search over candidate pairs or automorphisms occurs in verification.
GAP's exact permutation and subgroup algorithms are computational proof
infrastructure; they have not been replaced here by Lean proofs.
The complete checker, saved data and all required group identifications
remain inside the dependency boundary for any later comparison.

The other fourteen sporadic cases and the classical and exceptional
Lie families in the source remain separate obligations. This component
does not complete the original manuscript. The existing existence
results, the two printed Mathieu witnesses, the external group
classification and the geometric reconstruction theorem retain their
attribution; the new saved witnesses and this common checking argument
are the contribution of this component.

\begin{thebibliography}{9}
\bibitem{LL} D.~Leemans and M.~W.~Liebeck, \emph{Chiral polyhedra and
finite simple groups}, retained arXiv:1603.07713v3, Sections~1 and~6.
\bibitem{SW} E.~Schulte and A.~Ivi\'c Weiss, \emph{Chiral polytopes},
Applied Geometry and Discrete Mathematics, DIMACS Series 4 (1991), 493--516.
\bibitem{Atlas} J.~H.~Conway, R.~T.~Curtis, S.~P.~Norton, R.~A.~Parker
and R.~A.~Wilson, \emph{Atlas of Finite Groups}, Oxford University Press, 1985.
\bibitem{acquired} Human Compression project, \emph{Explicit sporadic
markings and their incidence certificates}, accepted N006 component;
exact files and admission recorded in the dependency ledger.
\end{thebibliography}
\end{document}
